\documentclass{oblivoir}
\usepackage{kotex}

\title{\textbf{2026 언어학과 교육환경총조사 수정안}}
\author{서울대학교 언어학과 서포터즈 2기 교육조장 이승헌}
\date{\today}

\setcounter{secnumdepth}{0}

\begin{document}
\maketitle
\tableofcontents

\newpage
\section{의제 3: 주전공 보호}
\subsection{2025 언어학과 교총조 결과}
\begin{enumerate}
    \item [3-1] 지금의 주전공 보호(강좌 분리 제도)가 주전공자의 학습권 보장이라는 목표를 달성하고 있다고 생각하십니까? (4점 척도) \\ 응답 결과: \textbf{\textit{2.3점}}
    \item [3-2] 3-1번과 같이 응답한 이유를 써 주세요. \\ a. 주전공자의 자유로운 분반 선택권이 보장되지 않으므로 \\ b. 주전공 분반(001)의 수강신청 난이도가 높아 전공 과목을 수강하지 못하는 경우가 있으므로 \\ c. (동일 과목) 분반별로 학습 내용에 차이가 존재하므로 \\ d. 기타
    \item [3-3] 신청일 분리 제도가 언어학과에서 시행되었을 때, 지금보다 주전공자의 학습권이 더 잘 보장될 것이라 생각하십니까? (4점 척도) \\ 응답 결과: \textbf{\textit{3.2점}} 
    \item [3-4] 3-3번과 같이 응답한 이유를 써 주세요. \\ a. 주전공생의 전공 과목 수강신청이 수월해질 것이므로 \\ b. 신청일 분리제도를 운영하고 있는 타과의 사례를 참고함 \\ c. 기타
    \item [3-5] 주전공 보호에 관하여 추가 의견이 있다면 자유롭게 써 주세요. \\ → 주전공생과 다전공생의 위계가 없었으면 좋겠음. 
\end{enumerate}

\subsection{2026 언어학과 교총조 설계}
\begin{enumerate}
    \item [3-1] 지금의 주전공 보호(강좌 분리 제도)가 주전공자의 학습권 보장이라는 목표를 달성하고 있지 않다는 의견이 주를 이루었음(2.3점 / 4점). 주된 이유가 무엇이라고 생각? \\ a. 주전공자의 자유로운 분반 선택권이 보장되지 않으므로 \\ b. 주전공 분반(001)의 수강신청 난이도가 높아 전공 과목을 수강하지 못하는 경우가 있으므로 \\ c. (동일 과목) 분반별로 학습 내용에 차이가 존재하므로 \\ d. 현행 제도가 주전공자의 학습권을 충분히 보장하고 있다고 생각 \\ e. 기타
    \item [3-2] 신청일 분리 제도가 언어학과에서 시행되었을 때, 지금보다 주전공자의 학습권이 더 잘 보장될 것이라는 의견이 주를 이루었음(3.2점/4점). 주된 이유가 무엇이라고 생각? \\ a. 주전공생의 전공 과목 수강신청이 수월해질 것이므로 \\ b. 신청일 분리제도를 운영하고 있는 타과의 사례를 참고함 \\ c. 주전공자의 학습권이 더 잘 보장될 것이라고 생각하지 않음 \\ d. 기타
    \item [3-3] 신청일 분리 제도는 주로 1일차(주전공생) - 2일차(다전공생) - 3일차(전체 학생) 순으로 수강신청 기회를 차등 부여함. 만약 신청일 분리 제도가 도입된다면, 이와 같이 기회를 차등 부여하는 것에 대해 동의? \\ a. 동의 \\ b. 비동의 \\ c. 모르겠음
    \item [3-4] (3-2에서 '비동의'에 답한 경우) 수강신청 기회가 구체적으로 어떻게 차등 부여되어야 한다고 생각하는지?
    \item [3-5] 주전공 보호에 관하여 추가 의견이 있다면 자유롭게 ㄱㄱ
\end{enumerate}

\newpage
\section{의제 7: 타 학과(부) 교과목 전공인정}
\subsection{2025 언어학과 교총조 결과}
\begin{enumerate}
    \item [7-1] 타 학과(부) 과목을 언어학과 전공교과목으로 인정받으신 적이 있습니까? \\ 응답 결과: \textbf{\textit{87\%}} (없음)
    \item [7-3] 타 학과(부)의 교과목을 수강하고 언어학과 전공교과목으로 인정받을 계획이 있으십니까? \\ 응답 결과: \textbf{\textit{69.6\%}} (없음)
    \item [7-5] 전공인정 신청을 하였으나 반려된 적이 있으십니까? \\ 응답 결과: \textbf{\textit{87\%}} (없음)
    \item [7-6] 현행 전공인정 방식에 대해 개편이 필요하다고 생각하십니까? \\ 응답 결과: \textbf{\textit{52.2\%}} (필요함), \textbf{\textit{47.8\%}} (잘 모르겠음)
    \item [7-7] (7-6에서 필요하다고 답한 경우) 전공인정 방식의 개편 방안 가운데, 동의하시는 것을 모두 선택해 주세요. \\ a. 인정되는 교과목 목록을 미리 정하여 공개 \textbf{\textit{92.3\%}} \\ b. 수강 전에 전공인정 신청 접수, 인정 가능 여부 통보 \textbf{\textit{69.2\%}}
\end{enumerate}

\subsection{2026 언어학과 교총조 설계}
\begin{enumerate}
    \item [7-1] TBA 
\end{enumerate}

\newpage
\section{의제 8: 기타}
2025 언어학과 교총조와 동일하게 구성. 끝.
\end{document}
